\chapter{问题重述与研究框架}

\section{问题重述}
“高精度参数确定问题”关注从有限观测数据中恢复动力系统参数。已知系统满足一组非线性常微分方程，观测数据通常包含测量误差，要求在拟合精度、参数可解释性与数值稳定性之间取得平衡。

\section{数据与材料说明}
压缩包中包含题目文档、\LaTeX{} 报告、导数表达及程序文件。本文重点利用以下材料：
\begin{itemize}
  \item \texttt{PredatorPreyTeX/PredatorPreyProblem.tex}：问题描述与原始推导框架；
  \item \texttt{ode56Order16.tex}, \texttt{RightFun56.tex}：高阶导数与右端项表达；
  \item C/Fortran/Matlab代码：数值积分与参数计算实现。
\end{itemize}

\section{总体技术路线}
本文采用四步法：
\begin{enumerate}
  \item 构建状态方程与观测方程；
  \item 建立最小二乘参数辨识目标；
  \item 设计高斯牛顿与阻尼牛顿迭代；
  \item 通过噪声扰动、初值扰动与多组实验验证鲁棒性。
\end{enumerate}
